% GENERATED FILE. Edit canonical lessons, then run npm run generate:books.
% canonical-source-hash: fnv1a64:0f9a423f94e762b1
% canonical-lessons: KA-C517-badti, KA-C517-pragati, KA-C517-mahatvakankse, KA-C517-sadhisu, KA-C517-galisu

\chapter{Getting Ahead}
\label{ch:ka-getting-ahead}
\hlchaptermodality{517}

\begin{chapteropening}
\textbf{By the end of this chapter:} \emph{I can say a promotion, progress, ambition, to achieve and to gain.}

Complete the last lesson of chapter 517: I can say a promotion, progress, ambition, to achieve and to gain.
\end{chapteropening}

\section[\texorpdfstring{\kn{ಬಡ್ತಿ} (baḍti)}{ಬಡ್ತಿ (baḍti)}]{\kn{ಬಡ್ತಿ} (baḍti) — a promotion}

\label{lesson:KA-C517-badti}



\begin{quote}
Before the new one: say the Kannada for an expert, then the Kannada for expertise.
\end{quote}

\subsection*{You'll want to know: \kn{ಬಡ್ತಿ}}
\textbf{\kn{ಬಡ್ತಿ}} — \emph{baḍti} — \textquotedblleft{}a promotion\textquotedblright{}.

\begin{quote}
\kn{ಈ} \kn{ವರ್ಷ} \kn{ನನ್ನ} \kn{ಅಕ್ಕನಿಗೆ} \kn{ಬಡ್ತಿ} \kn{ಸಿಕ್ಕಿತು}.
\end{quote}

\begin{quote}
\emph{ī varṣa nanna akkanige baḍti sikkitu.}
\end{quote}

\begin{quote}
This year my elder sister got a promotion.
\end{quote}

\kn{ಬಡ್ತಿ} \kn{ಕೊಡು}, to promote; \kn{ಬಡ್ತಿ} \kn{ಪಡೆ}, to be promoted.

\begin{cousinweb}
\textbf{\kn{ಬಡ್ತಿ}} (\emph{baḍti}) is usually traced to Hindustani \emph{baṛhtī}, \textquotedblleft{}growth, increase\textquotedblright{}; the path into Kannada is not settled.
\end{cousinweb}

\begin{grammarlens}[title={What you've built}]
One more word for describing what you lived and felt.
\end{grammarlens}

\subsection*{Your turn}
\begin{itemize}
\raggedright
  \item \emph{Say it:} \emph{baḍti}
  \item \emph{Say it:} \emph{baḍti}, once more
  \item \emph{Say it:} \emph{ī varṣa nanna akkanige baḍti sikkitu.} — This year my elder sister got a promotion.
  \item \emph{Recall:} say the Kannada for an expert, then the Kannada for expertise, then say \emph{baḍti} again
\end{itemize}

\begin{tcolorbox}[breakable,colback=teal!4,colframe=teal!35!black,title={Before you move on}]
What does \kn{ಬಡ್ತಿ} mean? (\textquotedblleft{}A promotion\textquotedblright{}.) Say it once more.
\end{tcolorbox}
\section[\texorpdfstring{\kn{ಪ್ರಗತಿ} (pragati)}{ಪ್ರಗತಿ (pragati)}]{\kn{ಪ್ರಗತಿ} (pragati) — progress}

\label{lesson:KA-C517-pragati}



\begin{quote}
Before the new one: say the Kannada for expertise, then the Kannada for a promotion.
\end{quote}

\subsection*{You'll want to know: \kn{ಪ್ರಗತಿ}}
\textbf{\kn{ಪ್ರಗತಿ}} — \emph{pragati} — \textquotedblleft{}progress\textquotedblright{}.

\begin{quote}
\kn{ಊರಿನಲ್ಲಿ} \kn{ರಸ್ತೆ} \kn{ಕೆಲಸ} \kn{ಪ್ರಗತಿಯಲ್ಲಿದೆ}.
\end{quote}

\begin{quote}
\emph{ūrinalli raste kelasa pragatiyallide.}
\end{quote}

\begin{quote}
Road work in the town is in progress.
\end{quote}

\kn{ಪ್ರಗತಿಪರ}, progressive.

\begin{cousinweb}
\textbf{\kn{ಪ್ರಗತಿ}} (\emph{pragati}) is built from Sanskrit \emph{pra-}, \textquotedblleft{}forward\textquotedblright{}, and \emph{gati}, \textquotedblleft{}going, progress\textquotedblright{} (Monier-Williams).
\end{cousinweb}

\begin{grammarlens}[title={What you've built}]
One more word for describing what you lived and felt.
\end{grammarlens}

\subsection*{Your turn}
\begin{itemize}
\raggedright
  \item \emph{Say it:} \emph{pragati}
  \item \emph{Say it:} \emph{pragati}, once more
  \item \emph{Say it:} \emph{ūrinalli raste kelasa pragatiyallide.} — Road work in the town is in progress.
  \item \emph{Recall:} say the Kannada for expertise, then the Kannada for a promotion, then say \emph{pragati} again
\end{itemize}

\begin{tcolorbox}[breakable,colback=teal!4,colframe=teal!35!black,title={Before you move on}]
What does \kn{ಪ್ರಗತಿ} mean? (\textquotedblleft{}Progress\textquotedblright{}.) Say it once more.
\end{tcolorbox}
\section[\texorpdfstring{\kn{ಮಹತ್ವಾಕಾಂಕ್ಷೆ} (mahatvākāṅkṣe)}{ಮಹತ್ವಾಕಾಂಕ್ಷೆ (mahatvākāṅkṣe)}]{\kn{ಮಹತ್ವಾಕಾಂಕ್ಷೆ} (mahatvākāṅkṣe) — ambition}

\label{lesson:KA-C517-mahatvakankse}



\begin{quote}
Before the new one: say the Kannada for a promotion, then the Kannada for progress.
\end{quote}

\subsection*{You'll want to know: \kn{ಮಹತ್ವಾಕಾಂಕ್ಷೆ}}
\textbf{\kn{ಮಹತ್ವಾಕಾಂಕ್ಷೆ}} — \emph{mahatvākāṅkṣe} — \textquotedblleft{}ambition\textquotedblright{}.

\begin{quote}
\kn{ಅವನಿಗೆ} \kn{ದೊಡ್ಡ} \kn{ವಿಜ್ಞಾನಿ} \kn{ಆಗುವ} \kn{ಮಹತ್ವಾಕಾಂಕ್ಷೆ} \kn{ಇದೆ}.
\end{quote}

\begin{quote}
\emph{avanige doḍḍa vijñāni āguva mahatvākāṅkṣe ide.}
\end{quote}

\begin{quote}
He has the ambition of becoming a great scientist.
\end{quote}

\kn{ಮಹತ್ವ}, greatness + \kn{ಆಕಾಂಕ್ಷೆ}, desire.

\begin{cousinweb}
\textbf{\kn{ಮಹತ್ವಾಕಾಂಕ್ಷೆ}} (\emph{mahatvākāṅkṣe}) is a learned compound built from Sanskrit \emph{mahattva}, \textquotedblleft{}greatness\textquotedblright{} (CDIAL 9933), and \emph{ākāṅkṣā}, \textquotedblleft{}desire, wish\textquotedblright{} (Monier-Williams).
\end{cousinweb}

\begin{grammarlens}[title={What you've built}]
One more word for describing what you lived and felt.
\end{grammarlens}

\subsection*{Your turn}
\begin{itemize}
\raggedright
  \item \emph{Say it:} \emph{mahatvākāṅkṣe}
  \item \emph{Say it:} \emph{mahatvākāṅkṣe}, once more
  \item \emph{Say it:} \emph{avanige doḍḍa vijñāni āguva mahatvākāṅkṣe ide.} — He has the ambition of becoming a great scientist.
  \item \emph{Recall:} say the Kannada for a promotion, then the Kannada for progress, then say \emph{mahatvākāṅkṣe} again
\end{itemize}

\begin{tcolorbox}[breakable,colback=teal!4,colframe=teal!35!black,title={Before you move on}]
What does \kn{ಮಹತ್ವಾಕಾಂಕ್ಷೆ} mean? (\textquotedblleft{}Ambition\textquotedblright{}.) Say it once more.
\end{tcolorbox}
\section[\texorpdfstring{\kn{ಸಾಧಿಸು} (sādhisu)}{ಸಾಧಿಸು (sādhisu)}]{\kn{ಸಾಧಿಸು} (sādhisu) — to achieve}

\label{lesson:KA-C517-sadhisu}



\begin{quote}
Before the new one: say the Kannada for progress, then the Kannada for ambition.
\end{quote}

\subsection*{You'll want to know: \kn{ಸಾಧಿಸು}}
\textbf{\kn{ಸಾಧಿಸು}} — \emph{sādhisu} — \textquotedblleft{}to achieve\textquotedblright{}.

\begin{quote}
\kn{ಕಷ್ಟಪಟ್ಟು} \kn{ಅವಳು} \kn{ತನ್ನ} \kn{ಕನಸನ್ನು} \kn{ಸಾಧಿಸಿದಳು}.
\end{quote}

\begin{quote}
\emph{kaṣṭapaṭṭu avaḷu tanna kanasannu sādhisidaḷu.}
\end{quote}

\begin{quote}
Through hard work she achieved her dream.
\end{quote}

Past (I …): \kn{ಸಾಧಿಸಿದೆ}, \emph{sādhiside}.

The noun is \kn{ಸಾಧನೆ}, an achievement.

\begin{cousinweb}
\textbf{\kn{ಸಾಧಿಸು}} (\emph{sādhisu}) is built on Sanskrit \emph{sādh-}, \textquotedblleft{}to accomplish\textquotedblright{}, whose noun \emph{sādhana} means \textquotedblleft{}furthering, completion\textquotedblright{} (CDIAL 13334), with the verb ending -\emph{isu}.
\end{cousinweb}

\begin{grammarlens}[title={What you've built}]
One more word for describing what you lived and felt.
\end{grammarlens}

\subsection*{Your turn}
\begin{itemize}
\raggedright
  \item \emph{Say it:} \emph{sādhisu}
  \item \emph{Say it:} \emph{sādhisu}, once more
  \item \emph{Say it:} \emph{kaṣṭapaṭṭu avaḷu tanna kanasannu sādhisidaḷu.} — Through hard work she achieved her dream.
  \item \emph{Recall:} say the Kannada for progress, then the Kannada for ambition, then say \emph{sādhisu} again
\end{itemize}

\begin{tcolorbox}[breakable,colback=teal!4,colframe=teal!35!black,title={Before you move on}]
What does \kn{ಸಾಧಿಸು} mean? (\textquotedblleft{}To achieve\textquotedblright{}.) Say it once more.
\end{tcolorbox}
\section[\texorpdfstring{\kn{ಗಳಿಸು} (gaḷisu)}{ಗಳಿಸು (gaḷisu)}]{\kn{ಗಳಿಸು} (gaḷisu) — to gain; to earn}

\label{lesson:KA-C517-galisu}



\begin{quote}
Before the new one: say the Kannada for ambition, then the Kannada for to achieve.
\end{quote}

\subsection*{You'll want to know: \kn{ಗಳಿಸು}}
\textbf{\kn{ಗಳಿಸು}} — \emph{gaḷisu} — \textquotedblleft{}to gain; to earn\textquotedblright{}.

\begin{quote}
\kn{ಅವನು} \kn{ಊರಿನಲ್ಲಿ} \kn{ಒಳ್ಳೆಯ} \kn{ಹೆಸರು} \kn{ಗಳಿಸಿದನು}.
\end{quote}

\begin{quote}
\emph{avanu ūrinalli oḷḷeya hesaru gaḷisidanu.}
\end{quote}

\begin{quote}
He earned a good name in the town.
\end{quote}

Past (I …): \kn{ಗಳಿಸಿದೆ}, \emph{gaḷiside}.

Money, fame or marks: \kn{ಅಂಕ} \kn{ಗಳಿಸು}, to score marks.

\begin{cousinweb}
\textbf{\kn{ಗಳಿಸು}} (\emph{gaḷisu}) is a native Kannada verb, \textquotedblleft{}to gain, to earn\textquotedblright{}; its origin is not settled.
\end{cousinweb}

\begin{grammarlens}[title={What you've built}]
One more word for describing what you lived and felt.
\end{grammarlens}

\subsection*{Your turn}
\begin{itemize}
\raggedright
  \item \emph{Say it:} \emph{gaḷisu}
  \item \emph{Say it:} \emph{gaḷisu}, once more
  \item \emph{Say it:} \emph{avanu ūrinalli oḷḷeya hesaru gaḷisidanu.} — He earned a good name in the town.
  \item \emph{Recall:} say the Kannada for ambition, then the Kannada for to achieve, then say \emph{gaḷisu} again
\end{itemize}

\begin{tcolorbox}[breakable,colback=teal!4,colframe=teal!35!black,title={Before you move on}]
What does \kn{ಗಳಿಸು} mean? (\textquotedblleft{}To gain; to earn\textquotedblright{}.) Say it once more.
\end{tcolorbox}
